% Vietnamese rendering of ai_statement.tex for the Thesis_v3 translation.
% The English source of truth remains ai_statement.tex at the repository root.
\newif\ifabstractunaided
\abstractunaidedtrue
\ifabstractunaided\else
  \GenericWarning{}{ai_statement: the abstract has not been rewritten unaided.
    Not submittable. Set \noexpand\abstractunaidedtrue in ai_statement.tex once it has}%
\fi

\cleardoublepage
\phantomsection
\addcontentsline{toc}{chapter}{Tuyên bố về việc dùng công cụ AI}
\chapter*{Tuyên bố về việc dùng công cụ AI}
\markboth{Tuyên bố về việc dùng công cụ AI}{Tuyên bố về việc dùng công cụ AI}

\noindent\textbf{Công cụ.} Claude Opus~5 của Anthropic và Cursor Grok~4.6 là các hệ thống AI
sinh được dùng khi soạn luận văn này. Cả hai chạy qua môi trường phát triển Cursor, từ tháng
6~2026 đến khi viết luận văn. Chính sách Đại học Helsinki cấm nêu một mô hình ngôn ngữ như một
tác giả. Không mô hình nào được trích như một nguồn, đó là lựa chọn của tác giả.

\noindent\textbf{Phạm vi hỗ trợ.} Các mô hình soạn thảo và sửa văn xuôi của luận văn, và đề xuất
thứ tự các lập luận. Chúng viết và tái cấu trúc cả hai bản dịch vụ Quản lý Ngữ cảnh UE. Chúng
cũng viết bộ sinh tải, danh mục tiêm lỗi, và nguyên mẫu FFI. Phân tích đứng sau các khoảng
bootstrap và các kiểm định hoán vị ghép cặp cũng là công việc của chúng, cũng như các nguồn hình,
build và LaTeX. Tóm tắt là bài kiểm tra trưởng thành của chương trình này. Chính sách đại học
không bao giờ cho phép một mô hình ngôn ngữ trong một bài như vậy.
\ifabstractunaided
Tóm tắt do tác giả viết không có hỗ trợ.
\else
Trong bản thảo này nó chưa được viết lại không có hỗ trợ, và tuyên bố này sẽ nói vậy cho đến khi
đã được viết lại.
\fi

\noindent\textbf{Những gì các mô hình không quyết.} Bài toán nghiên cứu, bốn câu hỏi nghiên cứu,
ba giả thuyết, và kiểm soát chính của nghiên cứu là của tác giả. Kiểm soát đó là một đặc tả trung
lập ngôn ngữ đã cố định và một mô hình tương tranh, nên ngôn ngữ vẫn là khoảng cách có kế hoạch
duy nhất giữa hai bản. Tác giả cũng chọn mang các lớp lỗi từ ngoài bảo đảm lúc biên dịch của
Rust, điều giữ RQ2 có thể bác bỏ. Mọi phép đo do tác giả chạy, trên các máy mà mỗi kết quả nêu.
Các đại lượng luận văn trích tới trang như macro được xây lại từ nhật ký thô. Cách đọc cũng là
của tác giả.

\noindent\textbf{Đầu ra được kiểm thế nào.} Mọi đầu ra của mô hình, văn xuôi và mã, đều được tác
giả đọc trước khi vào luận văn. Mọi mục thư mục được đối chiếu với Crossref, DataCite, hoặc cơ
quan phát hành. Một cuộc kiểm rồi chấm mức mỗi tuyên bố truy được. Bậc thứ nhất là chính trang
hoặc mệnh đề. Bậc thứ hai là hồ sơ nhà xuất bản nộp mang tuyên bố. Bậc thứ ba là siêu dữ liệu đã
xác nhận sau tường phí. Một mục mà đoạn chống đỡ chưa bao giờ được mở thì được liệt kê, không
được hấp thụ. Cuộc kiểm tìm bốn đoạn phóng đại hoặc đặt sai nguồn. Một đoạn là một trích dẫn vắng
trong công trình được gán. Cả bốn đã được sửa. Mã được kiểm đối với một bộ tuân thủ mà bản ghi
được diff từng dòng. Nó cũng được kiểm đối với một danh mục lỗi chốt mỗi kết cục kỳ vọng trước
khi lỗi được tiêm.

\noindent\textbf{Trách nhiệm.} Tác giả đã rà toàn bộ luận văn này và chịu trách nhiệm về tất cả.
Nơi nó sai, lỗi là của tác giả.
